\section{Experimental setup}
\label{sec:setup}

\subsection{Hardware, software and decoding}
\label{sec:setup-hardware}

All local experiments run on one laptop: an AMD Ryzen~7 8845HS CPU with
16\,GB of system memory and an NVIDIA GeForce RTX~4060 Laptop GPU with 8\,GB
of memory, driver 616.56 and an enforced board power limit of 115\,W. The
driver was updated to 616.92 between the main labelling run and the later
stages; the noise stage, which resamples prompts and cells the main run
measured, shows energy per output token under 616.92 at 0.89 to 1.01 times
that under 616.56. Every cell of the grid was measured under 616.56, so no
comparison between cells crosses the change. Models are served by Ollama \cite{ollama} from GGUF
weights quantised to Q4\_K\_M \cite{gerganov2023llamacpp}. Each model is
addressed by an explicit quantisation tag, and its weight digest, parameter
count, context length and chat template are recorded in a manifest before any
generation, so a silently updated tag cannot change the models between runs.

The 8\,GB memory limit means one model is loaded at a time, so runs are
organised \emph{model-major}: each model is loaded once, all prompts run in
both modes, and the model is unloaded before the next. Within a model,
prompts are processed in batches of 20, each batch running its two modes as
separate measurement windows with the leading mode alternating, so that slow
drift in temperature or background load cannot line up with one mode. The
harness waits for 30\,s of continuous GPU idleness before loading a model or
repeating a window, and checks that the whole model is placed on the GPU.

Greedy decoding is discouraged for Qwen3's thinking mode, so each mode uses
its recommended sampling settings \cite{qwen2025modelcard}: temperature 0.6,
top-$p$ 0.95 and top-$k$ 20 for thinking, and 0.7, 0.8 and 20 for
non-thinking. Every generation receives its own seed, derived
deterministically from the prompt, cell, sample index and run variant, and
stored with the output. The labelling runs draw one sample per prompt and
cell ($k = 1$), so a noise subset of 25 prompts per family receives two
further samples in every cell to measure how often a label would change on
resampling; three samples throughout were the original plan, but one sample
across the eight cells takes about 112\,s per prompt, and tripling that would
have tripled a run that already takes days.

Generation length is capped, and an answer that hits the cap counts as a
failure, because a truncated answer is of no use in deployment: 4{,}096 new
tokens in thinking mode, trace included, and 2{,}048 without. The caps matter
more than we expected, so their effect is measured rather than assumed. A
sample of capped generations, up to 40 per size and mode and spread evenly
over families, is rerun with an 8{,}192-token cap and a 10{,}240-token
context (9{,}216 at 8B, where the larger context no longer fits in GPU
memory), and we report how often the longer cap turns a failure into a
success and how often that moves $\cstar(q)$. Each server session runs one
thinking and one non-thinking warm-up generation, which are discarded.

\subsection{Corpus, labelling and grading}
\label{sec:setup-labels}

The corpus covers four task families of 400 prompts each. Three are graded
programmatically: mathematics, from GSM8K \cite{cobbe2021gsm8k} (200) and
MATH-500 \cite{hendrycks2021math,lightman2024verify} (200); code, from
HumanEval \cite{chen2021humaneval} (154) and MBPP \cite{austin2021mbpp} (246);
and multiple choice, from MMLU \cite{hendrycks2021mmlu} (200) and
ARC-Challenge \cite{clark2018arc} (200). The fourth, open-ended requests from
Dolly-15k \cite{conover2023dolly} (400), is the traffic that motivates the
problem and needs a judge (Section~\ref{sec:setup-judge}).

Prompts are drawn with a fixed, recorded seed. Each source is walked in
round-robin order over strata of an a-priori difficulty proxy where one
exists---MATH level, MMLU subject and MBPP task band---so that any prefix of
the order is balanced across strata; a 100-prompt pilot per family is
therefore nested in the full corpus, and its first 25 prompts form the noise
subset. Near-duplicates are removed by hashing normalised prompt text across
sources as well as within them, and prompts longer than 2{,}048 tokens are
excluded so that context length does not vary with the size axis.

Every prompt is run through all eight cells, recording the full success
vector $s(q, c_1), \dots, s(q, c_8)$ with the token counts, latency and
energy of every cell. Routers are trained on that vector, which carries more
information than $\cstar(q)$ and is what allows $\tau$ to be swept at test
time. Success is decided per family: for mathematics the final answer is
extracted and normalised before comparison with the reference; for code the
generated function is executed against the dataset's unit tests in a
subprocess with a timeout and no network access, and succeeds only if every
test passes; for multiple choice the chosen option is parsed, and an
unparseable answer is a failure rather than a reason to resample. After
labelling, the corpus is split 60/20/20 into training, validation and test
sets by prompt, stratified by family and by $\cstar(q)$ computed in tokens,
so that every cost unit is evaluated on the same split; the test split is
used once, at the end.

\subsection{Judging open-ended answers}
\label{sec:setup-judge}

A judge drawn from the routed pool would grade its own family, so the judge
is Llama~3.1 8B Instruct \cite{grattafiori2024llama3} at Q4\_K\_M, run only
after labelling has finished. It sees the request, the human-written Dolly
response as a reference and the visible answer, and scores the answer from 1
to 5 on a fixed rubric, with 4 or more a success; for open requests such as
brainstorming the rubric treats the reference as one acceptable answer among
many. Answers are scored one at a time in a seeded random order, to limit the
position and verbosity biases documented for model judges
\cite{zheng2023judging}, with greedy decoding constrained to a JSON reply.

The judge is validated from two sides before its labels are used. About 200
judged answers, stratified by category and verdict, are re-scored with the
same rubric by \texttt{gpt-5-mini} through a hosted API, a reference above
the local judge in capability and outside its family; and the first author
scores about 100 of them blind to model, mode and judge score. Agreement is
Cohen's $\kappa$ \cite{cohen1960kappa} on the success decision and
quadratic-weighted $\kappa$ \cite{cohen1968weighted} on the scores. The rule
for using the labels was fixed before any judge output existed: at
$\kappa \geq 0.60$ the open-ended family is reported with the others; from
0.40 to 0.60 it carries an explicit caveat; below 0.40 it moves to a
secondary analysis, these being the conventional boundaries of moderate and
substantial agreement \cite{landis1977measurement}. The judge scored all
3{,}583 eligible answers, reaching $\kappa = 0.42$ [0.23, 0.59] against the
first author's blind scores and 0.40 [0.27, 0.52] against
\texttt{gpt-5-mini}, which agrees with the first author at 0.74 [0.60, 0.86].
The family therefore carries the caveat and is reported separately.

\subsection{Baselines, metrics and statistics}
\label{sec:setup-stats}

Routers are compared with the eight static policies that always use one cell,
uniform random assignment, the oracle that always picks $\cstar(q)$, the four
manual policies of Table~\ref{tab:policies} under both detection regimes, and
the cascade of Section~\ref{sec:net}. We report cost--quality frontiers and
the normalised area beneath them, the tax and its components from
\eqref{eq:decomp}, net savings from \eqref{eq:netsaving}, cost per solved
query in each unit, the over- and under-routing rates with their separate
cost consequences, and router overhead per decision.

Every headline figure carries a 95\% percentile bootstrap confidence interval
from 10{,}000 resamples over prompts \cite{efron1993bootstrap}. All policies
and routers evaluated on the same prompts share one matrix of resampled
indices, so every comparison is paired by prompt and its interval is taken
over resampled differences. Two-sided $p$-values are obtained by inverting
the percentile interval, and the Holm--Bonferroni procedure
\cite{holm1979simple} corrects within each family of comparisons that asks
one question at a fixed cost unit, recovery rule and detection regime. For
routers, $\tau$ is fixed on the validation split before the test split is
resampled, so the intervals exclude the uncertainty of that choice. Each
trainable router is trained with three seeds, with intervals reported for the
first and the spread across seeds beside them.

\subsection{Validation gates}
\label{sec:setup-gates}

Two assumptions were tested before any labelling, and both gates are reported
in full in the supplementary material (Tables S2 and S3).

\emph{Gate 1, power telemetry.} NVML does not expose board power on every
laptop GPU. Recording 10\,s of idle, 8\,s of half-precision matrix
multiplication at about 90\% utilisation and 8\,s of cool-down, sampled power
behaved as expected: low at idle, following the step load within about a
second, and never exceeding the enforced limit. The cumulative energy counter
did not. It is reported as supported and increases monotonically, so a
routine check accepts it, yet the power it implies is 88\,W with the GPU
idle, 128\,W under load, above the enforced limit, and 257\,W during
cool-down, higher than under load. A units error would scale all three
windows by one factor and keep their order; these values do not follow real
power draw at all. We therefore integrate sampled power only, and the lesson
for other studies on laptop GPUs is that a supported, monotonic counter
should still be cross-checked against sampled power and the enforced limit.

\emph{Gate 2, separation of thinking modes.} The grid assumes a model's two
modes are different operating points, which Wang et al.\ found only partly
true of hybrid-thinking models \cite{wang2025demystifying}. On 60 prompts per
size, run once per mode under a common 4{,}096-token cap, three criteria
fixed in advance were all met at all four sizes: an area under the ROC curve
of at least 0.90 for thinking outputs being longer (observed 0.928 to 0.964),
at most 5\% of non-thinking responses containing a thinking trace or tag
(observed zero), and reasoning-marker density in non-thinking answers no more
than half that in thinking traces (observed 3--5\%). Marker counting is a
lexical proxy, and 60 prompts per size cannot rule out leakage on task types
we did not sample. The gate also showed that thinking dominates the compute
budget, taking 88\% of the time of one sample across the eight cells, and
that between 21\% and 32\% of thinking responses hit the cap, which is why
the cap-sensitivity reruns of Section~\ref{sec:setup-hardware} are part of
the design.

\subsection{Public-benchmark and hosted-frontier arms}
\label{sec:setup-arms}

The local pool stops at 8B parameters, so two further arms test whether the
analysis carries beyond it. The first applies the same policy and tax
computations, without local inference, to the published per-model outcomes
and costs of LLMRouterBench \cite{li2026llmrouterbench}, whose records store a
score, prompt and completion token counts and the provider's dollar charge
for one model on one query---which is what the resubmission term needs, since
a re-asked query pays its prompt tokens again at the stronger model's price.
We take the thirteen models of its performance--cost setting and the nine
datasets every one of them was run on, keeping per dataset the largest subset
all thirteen share, which leaves 12{,}168 queries. A pool of thirteen models
has no natural order, so the ladder is the cost--accuracy frontier in cost
order, whose top rung plays the part $c_8$ plays locally; P2 has no analogue
without a mode axis, and the thinking and non-thinking releases of
Qwen3-235B-A22B are compared separately in its place. Because the frontier is
a subset of the pool, a policy climbing it can fail on a query some
off-ladder model answers, so the tax is read alongside the success rate.
Intervals here use 2{,}000 resamples, the prompt set being ten times the
local corpus.

The second arm restores the mode axis at frontier scale. Three hosted models,
\texttt{gpt-5-nano}, \texttt{gpt-5-mini} and \texttt{gpt-5}, are crossed with
two reasoning efforts, \texttt{minimal} and \texttt{high}, giving six cells
$h_1 \dots h_6$ indexed as $c_1 \dots c_8$ are, priced at list prices
published on 12 September 2026 and halved by the batch discount. Prompts come
from those the local grid has already labelled, so every hosted $\cstar$ is
paired with a local one. Each request carries exactly the prompt text the
local pool saw, and the output caps mirror the local protocol; on a reasoning
model the cap counts reasoning tokens, so a truncated response is a failure
exactly as a local cap hit is. That rule falls hardest on high-effort cells,
so each cell is reported twice: the headline accuracy, in which truncation is
wrong, and the accuracy over generations that stopped on their own. Three
limits are worth stating. The hosted models reject the temperature and
nucleus parameters the local protocol sets, so the arm matches the prompt and
the cap, not the decoding. Reasoning effort is not the local binary switch,
though at \texttt{minimal} each model reported exactly zero reasoning tokens
across all 120 pilot responses, so the cheap mode really is reasoning off.
And the provider returns the number of reasoning tokens but not the trace,
and no energy, so this arm reports dollars and tokens only. A prompt enters
the analysis only when all six cells answered it, which two of 827 prompts
failed through content-filter refusals, leaving 825.

\subsection{Reproducibility}
\label{sec:setup-repro}

Each experiment is driven by a configuration file, and no experimental
constant lives in code. Seeds are set and stored for every generation and
training run, and each run writes a manifest with the code commit,
configuration, package versions, driver version, GPU model and model digests.
Raw generations are stored with their token counts, timings, seeds and
energy. The hosted arm adds a spend ledger recording every batch with the
worst case it was approved against and the cost finally settled from the
provider's own token counts, so each dollar figure can be audited line by
line; the total agrees to the cent with the provider's billed figure. Hosted
models are still not reproducible artefacts in the way a pinned local
checkpoint is, which is why the local grid carries the primary claims. The
analysis code is archived as \cite{mondal2026arbiter}.
